\section{Universally injective module maps}
\label{section-universally-injective}

\noindent
Next we discuss universally injective module maps, which
are in a sense complementary to flat modules (see Lemma
\ref{lemma-flat-universally-injective}).  We follow Lazard's thesis
\cite{Autour};  also see \cite{Lam}.

\begin{definition}
\label{definition-universally-injective}
Let $f: M \to N$ be a map of $R$-modules.  Then $f$ is called
{\it universally injective} if for every $R$-module $Q$, the map $f
\otimes_R \text{id}_Q: M \otimes_R Q \to N \otimes_R Q$
is injective.  A sequence $0 \to M_1 \to M_2 \to M_3
\to 0$ of $R$-modules is called {\it universally exact} if it is exact
and $M_1 \to M_2$ is universally injective.
\end{definition}

\begin{example}
\label{example-universally-exact}
Examples of universally exact sequences.
\begin{enumerate}
\item A split short exact sequence is universally exact since tensoring
commutes with taking direct sums.
\item The colimit of a directed system of universally exact sequences is
universally exact.  This follows from the fact that taking directed colimits is
exact and that tensoring commutes with taking colimits.  In particular the
colimit of a directed system of split exact sequences is universally exact.  We
will see below that, conversely, any universally exact sequence arises in this
way.
\end{enumerate}
\end{example}

\noindent
Next we give a list of criteria for a short exact sequence to be universally
exact.  They are analogues of criteria for flatness given above.  Parts (3)-(6)
below correspond, respectively, to the criteria for flatness given in
Lemmas \ref{lemma-flat-eq}, \ref{lemma-flat-factors-free},
\ref{lemma-flat-surjective-hom}, and
Theorem \ref{theorem-lazard}.

\begin{theorem}
\label{theorem-universally-exact-criteria}
Let
$$
0 \to M_1 \xrightarrow{f_1} M_2 \xrightarrow{f_2} M_3 \to 0
$$
be an exact sequence of $R$-modules.  The following are equivalent:
\begin{enumerate}
\item The sequence $0 \to M_1 \to M_2 \to M_3
\to 0$ is universally exact.
\item For every finitely presented $R$-module $Q$, the sequence
$$
0 \to M_1 \otimes_R Q \to M_2 \otimes_R Q \to
M_3 \otimes_R Q \to 0
$$
is exact.
\item Given elements $x_i \in M_1$ $(i = 1, \ldots, n)$, $y_j \in M_2$ $(j = 1,
\ldots, m)$, and $a_{ij} \in R$ $(i = 1, \ldots, n, j = 1, \ldots, m)$ such that
for all $i$
$$
f_1(x_i) = \sum\nolimits_j a_{ij} y_j,
$$
there exists $z_j \in M_1$ $(j =1, \ldots, m)$ such that for all $i$,
$$
x_i = \sum\nolimits_j a_{ij} z_j .
$$
\item Given a commutative diagram of $R$-module maps
$$
\xymatrix{
R^n \ar[r] \ar[d] &  R^m \ar[d] \\
M_1 \ar[r]^{f_1}        &  M_2
}
$$
where $m$ and $n$ are integers, there exists a map $R^m \to M_1$ making
the top triangle commute.
\item For every finitely presented $R$-module $P$, the $R$-module
map $\Hom_R(P, M_2) \to \Hom_R(P, M_3)$ is surjective.
\item The sequence $0 \to M_1 \to M_2 \to M_3
\to 0$ is the colimit of a directed system of split exact sequences of
the form
$$
0 \to M_{1} \to M_{2, i} \to M_{3, i} \to 0
$$
where the $M_{3, i}$ are finitely presented.
\end{enumerate}
\end{theorem}

\begin{proof}
Obviously (1) implies (2).

\medskip\noindent
Next we show (2) implies (3).  Let $f_1(x_i) = \sum_j a_{ij} y_j$ be relations
as in (3).  Let $(d_j)$ be a basis for $R^m$, $(e_i)$ a basis for $R^n$, and
$R^m \to R^n$ the map given by $d_j \mapsto \sum_i a_{ij} e_i$.
Let $Q$ be the cokernel of $R^m \to R^n$.  Then tensoring
$R^m \to R^n \to Q \to 0$ by the map $f_1: M_1 \to M_2$, we get a
commutative diagram
$$
\xymatrix{
M_1^{\oplus m} \ar[r] \ar[d] & M_1^{\oplus n} \ar[r] \ar[d] & M_1 \otimes_R Q
\ar[r] \ar[d] & 0 \\
M_2^{\oplus m} \ar[r] & M_2^{\oplus n} \ar[r] & M_2 \otimes_R Q \ar[r] & 0
}
$$
where $M_1^{\oplus m} \to M_1^{\oplus n}$ is given by
$$
(z_1, \ldots, z_m) \mapsto
(\sum\nolimits_j a_{1j} z_j, \ldots, \sum\nolimits_j a_{nj} z_j),
$$
and $M_2^{\oplus m} \to M_2^{\oplus n}$ is given similarly.  We want to
show $x = (x_1, \ldots, x_n) \in M_1^{\oplus n}$ is in the image of $M_1^{\oplus
m} \to M_1^{\oplus n}$.  By (2) the map $M_1 \otimes Q \to M_2
\otimes Q$ is injective, hence by exactness of the top row it is enough to show
$x$ maps to $0$ in $M_2 \otimes Q$, and so by exactness of the bottom row it is
enough to show the image of $x$ in $M_2^{\oplus n}$ is in the image of
$M_2^{\oplus m} \to M_2^{\oplus n}$.  This is true by assumption.

\medskip\noindent
Condition (4) is just a translation of (3) into diagram form.

\medskip\noindent
Next we show (4) implies (5). Let $\varphi : P \to M_3$ be a map from a
finitely presented $R$-module $P$.  We must show that $\varphi$ lifts to a map
$P \to M_2$.  Choose a presentation of $P$,
$$
R^n \xrightarrow{g_1} R^m \xrightarrow{g_2} P \to 0.
$$
Using freeness of $R^n$ and $R^m$, we can construct $h_2: R^m \to M_2$
and then $h_1: R^n \to M_1$ such that the following diagram commutes
$$
\xymatrix{
         & R^n \ar[r]^{g_1} \ar[d]^{h_1} & R^m \ar[r]^{g_2} \ar[d]^{h_2} & P
\ar[r] \ar[d]^{\varphi} & 0 \\
0 \ar[r] & M_1 \ar[r]^{f_1} & M_2 \ar[r]^{f_2} & M_3 \ar[r] & 0 .
}
$$
By (4) there is a map $k_1: R^m \to M_1$ such that $k_1 \circ g_1 =
h_1$.  Now define $h'_2: R^m \to M_2$ by $h_2' = h_2 - f_1 \circ k_1$.
Then
$$
h'_2 \circ g_1 = h_2 \circ g_1 - f_1 \circ k_1 \circ g_1 =
h_2 \circ g_1 - f_1 \circ h_1 = 0 .
$$
Hence by passing to the quotient $h'_2$ defines a map $\varphi': P \to
M_2$ such that $\varphi' \circ g_2 = h_2'$.  In a diagram, we have
$$
\xymatrix{
R^m \ar[r]^{g_2} \ar[d]_{h'_2} & P \ar[d]^{\varphi} \ar[dl]_{\varphi'} \\
M_2 \ar[r]^{f_2} & M_3.
}
$$
where the top triangle commutes.  We claim that $\varphi'$ is the desired lift,
i.e.\ that $f_2 \circ \varphi' = \varphi$.  From the definitions we have
$$
f_2 \circ \varphi' \circ g_2 = f_2 \circ h'_2 =
f_2 \circ h_2 - f_2 \circ f_1 \circ k_1 = f_2 \circ h_2 =
\varphi \circ g_2.
$$
Since $g_2$ is surjective, this finishes the proof.

\medskip\noindent
Now we show (5) implies (6). Write $M_{3}$ as the colimit of a directed system
of finitely presented modules $M_{3, i}$, see
Lemma \ref{lemma-module-colimit-fp}. Let $M_{2, i}$ be the fiber product
of $M_{3, i}$ and $M_{2}$ over $M_{3}$---by definition this is the submodule of
$M_2 \times M_{3, i}$ consisting of elements whose two projections onto $M_3$
are equal. Let $M_{1, i}$ be the kernel of the projection
$M_{2, i} \to M_{3, i}$. Then we have a directed system of exact sequences
$$
0 \to M_{1, i} \to M_{2, i} \to M_{3, i} \to 0,
$$
and for each $i$ a map of exact sequences
$$
\xymatrix{
0  \ar[r] & M_{1, i} \ar[d] \ar[r]  &  M_{2, i}  \ar[r] \ar[d] & M_{3, i} \ar[d]
\ar[r] & 0 \\
0  \ar[r] & M_{1} \ar[r]  &  M_{2}  \ar[r] & M_{3} \ar[r] & 0
}
$$
compatible with the directed system.  From the definition of the fiber product
$M_{2, i}$, it follows that the map $M_{1, i} \to M_1$ is an isomorphism.
 By (5) there is a map $M_{3, i} \to M_{2}$ lifting $M_{3, i} \to
M_3$, and by the universal property of the fiber product this gives rise to a
section of $M_{2, i} \to M_{3, i}$.  Hence the sequences
$$
0 \to M_{1, i} \to M_{2, i} \to M_{3, i} \to 0
$$
split.  Passing to the colimit, we have a commutative diagram
$$
\xymatrix{
0  \ar[r] & \colim M_{1, i} \ar[d]^{\cong} \ar[r]  &  \colim M_{2, i}
 \ar[r]
\ar[d] & \colim M_{3, i} \ar[d]^{\cong} \ar[r] & 0 \\
0  \ar[r] & M_{1} \ar[r]  &  M_{2}  \ar[r] & M_{3} \ar[r] & 0
}
$$
with exact rows and outer vertical maps isomorphisms.  Hence $\colim
M_{2, i}
\to M_2$ is also an isomorphism and (6) holds.

\medskip\noindent
Condition (6) implies (1) by
Example \ref{example-universally-exact} (2).
\end{proof}

\noindent
The previous theorem shows that a universally exact sequence is always a
colimit of split short exact sequences.  If the cokernel of a universally
injective map is finitely presented, then in fact the map itself splits:

\begin{lemma}
\label{lemma-universally-exact-split}
Let
$$
0 \to M_1 \to M_2 \to M_3 \to 0
$$
be an exact sequence of $R$-modules.  Suppose $M_3$ is of finite presentation.
Then
$$
0 \to M_1 \to M_2 \to M_3 \to 0
$$
is universally exact if and only if it is split.
\end{lemma}

\begin{proof}
A split short exact sequence is always universally exact, see
Example \ref{example-universally-exact}.
Conversely, if the sequence is universally exact, then by
Theorem \ref{theorem-universally-exact-criteria} (5)
applied to $P = M_3$, the map $M_2 \to M_3$ admits a section.
\end{proof}

\noindent
The following lemma shows how universally injective maps are complementary to
flat modules.

\begin{lemma}
\label{lemma-flat-universally-injective}
Let $M$ be an $R$-module.  Then $M$ is flat if and only if any exact sequence
of $R$-modules
$$
0 \to M_1 \to M_2 \to M \to 0
$$
is universally exact.
\end{lemma}

\begin{proof}
This follows from Lemma \ref{lemma-flat-surjective-hom} and
Theorem \ref{theorem-universally-exact-criteria} (5).
\end{proof}

\begin{example}
\label{example-universally-exact-non-split-non-flat}
Non-split and non-flat universally exact sequences.
\begin{enumerate}
\item In spite of
Lemma \ref{lemma-universally-exact-split},
it is possible to  have a short exact sequence of $R$-modules
$$
0 \to M_1 \to M_2 \to M_3 \to 0
$$
that is universally exact but non-split.  For instance, take $R = \mathbf{Z}$,
let $M_1 = \bigoplus_{n=1}^{\infty} \mathbf{Z}$, let $M_{2} = \prod_{n =
1}^{\infty} \mathbf{Z}$, and let $M_{3}$ be the cokernel of the inclusion $M_1
\to M_2$.  Then $M_1, M_2, M_3$ are all flat since they are torsion-free
(More on Algebra, Lemma \ref{more-algebra-lemma-dedekind-torsion-free-flat}),
so by Lemma \ref{lemma-flat-universally-injective},
$$
0 \to M_1 \to M_2 \to M_3 \to 0
$$
is universally exact.  However there can be no section $s: M_3 \to
M_2$.  In fact, if $x$ is the image of $(2, 2^2, 2^3, \ldots) \in M_2$ in $M_3$,
then any module map $s: M_3 \to M_2$ must kill $x$.  This is because $x
\in 2^n M_3$ for any $n \geq 1$, hence $s(x)$ is divisible by $2^n$ for all $n
\geq 1$ and so must be $0$.
\item In spite of Lemma \ref{lemma-flat-universally-injective}, it is possible
to have a short exact sequence of $R$-modules
$$
0 \to M_1 \to M_2 \to M_3 \to 0
$$
that is universally exact but with $M_1, M_2, M_3$ all non-flat.  In fact if $M$
is any non-flat module, just take the split exact sequence
$$
0 \to M \to M \oplus M \to M \to 0.
$$
For instance over $R = \mathbf{Z}$, take $M$ to be any torsion module.
\item Taking the direct sum of an exact sequence as in (1) with one as in (2),
we get a short exact sequence of $R$-modules
$$
0 \to M_1 \to M_2 \to M_3 \to 0
$$
that is universally exact, non-split, and such that
$M_1, M_2, M_3$ are all non-flat.
\end{enumerate}
\end{example}

\begin{lemma}
\label{lemma-ui-flat-domain}
Let $0 \to M_1 \to M_2 \to M_3 \to 0$ be a universally exact sequence
of $R$-modules, and suppose $M_2$ is flat.
Then $M_1$ and $M_3$ are flat.
\end{lemma}

\begin{proof}
Let $0 \to N \to N' \to N'' \to 0$ be a short exact sequence of
$R$-modules. Consider the commutative diagram
$$
\xymatrix{
M_1 \otimes_R N \ar[r] \ar[d] &
M_2 \otimes_R N \ar[r] \ar[d] &
M_3 \otimes_R N \ar[d] \\
M_1 \otimes_R N' \ar[r] \ar[d] &
M_2 \otimes_R N' \ar[r] \ar[d] &
M_3 \otimes_R N' \ar[d] \\
M_1 \otimes_R N'' \ar[r] &
M_2 \otimes_R N'' \ar[r] &
M_3 \otimes_R N''
}
$$
(we have dropped the $0$'s on the boundary).
By assumption the rows give short exact sequences and the arrow
$M_2 \otimes N \to M_2 \otimes N'$ is injective. Clearly this implies
that $M_1 \otimes N \to M_1 \otimes N'$ is injective and we see that $M_1$
is flat. In particular the left and middle columns give rise to short
exact sequences. It follows from a diagram chase that the arrow
$M_3 \otimes N \to M_3 \otimes N'$ is injective. Hence $M_3$ is flat.
\end{proof}

\begin{lemma}
\label{lemma-universally-injective-tensor}
Let $R$ be a ring.
Let $M \to M'$ be a universally injective $R$-module map.
Then for any $R$-module $N$ the map $M \otimes_R N \to M' \otimes_R N$
is universally injective.
\end{lemma}

\begin{proof}
Omitted.
\end{proof}

\begin{lemma}
\label{lemma-composition-universally-injective}
Let $R$ be a ring. A composition of universally injective
$R$-module maps is universally injective.
\end{lemma}

\begin{proof}
Omitted.
\end{proof}

\begin{lemma}
\label{lemma-universally-injective-permanence}
Let $R$ be a ring. Let $M \to M'$ and $M' \to M''$ be $R$-module maps.
If their composition $M \to M''$ is universally injective, then
$M \to M'$ is universally injective.
\end{lemma}

\begin{proof}
Omitted.
\end{proof}

\begin{lemma}
\label{lemma-universally-injective-check-stalks}
Let $R \to S$ be a ring map.
Let $M \to M'$ be a map of $S$-modules.
The following are equivalent
\begin{enumerate}
\item $M \to M'$ is universally injective as a map of $R$-modules,
\item for each prime $\mathfrak q$ of $S$ the map
$M_{\mathfrak q} \to M'_{\mathfrak q}$ is universally injective
as a map of $R$-modules,
\item for each maximal ideal $\mathfrak m$ of $S$ the map
$M_{\mathfrak m} \to M'_{\mathfrak m}$ is universally injective
as a map of $R$-modules,
\item for each prime $\mathfrak q$ of $S$ the map
$M_{\mathfrak q} \to M'_{\mathfrak q}$ is universally injective
as a map of $R_{\mathfrak p}$-modules, where $\mathfrak p$ is the
inverse image of $\mathfrak q$ in $R$, and
\item for each maximal ideal $\mathfrak m$ of $S$ the map
$M_{\mathfrak m} \to M'_{\mathfrak m}$ is universally injective
as a map of $R_{\mathfrak p}$-modules, where $\mathfrak p$ is the
inverse image of $\mathfrak m$ in $R$.
\end{enumerate}
\end{lemma}

\begin{proof}
Let $N$ be an $R$-module. Let $\mathfrak q$ be a prime of $S$ lying over
the prime $\mathfrak p$ of $R$. Then we have
$$
(M \otimes_R N)_{\mathfrak q} =
M_{\mathfrak q} \otimes_R N =
M_{\mathfrak q} \otimes_{R_{\mathfrak p}} N_{\mathfrak p}.
$$
Moreover, the same thing holds for $M'$ and localization is exact.
Also, if $N$ is an $R_{\mathfrak p}$-module, then $N_{\mathfrak p} = N$.
Using this the equivalences can be proved in a straightforward manner.

\medskip\noindent
For example, suppose that (5) holds. Let
$K = \Ker(M \otimes_R N \to M' \otimes_R N)$. By the remarks
above we see that $K_{\mathfrak m} = 0$ for each maximal ideal $\mathfrak m$
of $S$. Hence $K = 0$ by
Lemma \ref{lemma-characterize-zero-local}.
Thus (1) holds. Conversely, suppose that (1) holds. Take any
$\mathfrak q \subset S$ lying over $\mathfrak p \subset R$.
Take any module $N$ over $R_{\mathfrak p}$. Then
by assumption $\Ker(M \otimes_R N \to M' \otimes_R N) = 0$.
Hence by the formulae above and the fact that $N = N_{\mathfrak p}$
we see that
$\Ker(M_{\mathfrak q} \otimes_{R_{\mathfrak p}} N \to
M'_{\mathfrak q} \otimes_{R_{\mathfrak p}} N) = 0$. In other words
(4) holds. Of course (4) $\Rightarrow$ (5) is immediate. Hence
(1), (4) and (5) are all equivalent.
We omit the proof of the other equivalences.
\end{proof}

\begin{lemma}
\label{lemma-universally-injective-localize}
Let $\varphi : A \to B$ be a ring map. Let $S \subset A$ and
$S' \subset B$ be multiplicative subsets such that $\varphi(S) \subset S'$.
Let $M \to M'$ be a map of $B$-modules.
\begin{enumerate}
\item If $M \to M'$ is universally injective as a map of $A$-modules,
then $(S')^{-1}M \to (S')^{-1}M'$ is universally injective as a map of
$A$-modules and as a map of $S^{-1}A$-modules.
\item If $M$ and $M'$ are $(S')^{-1}B$-modules, then $M \to M'$
is universally injective as a map of $A$-modules if and only if
it is universally injective as a map of $S^{-1}A$-modules.
\end{enumerate}
\end{lemma}

\begin{proof}
You can prove this using
Lemma \ref{lemma-universally-injective-check-stalks}
but you can also prove it directly as follows.
Assume $M \to M'$ is $A$-universally injective.
Let $Q$ be an $A$-module. Then $Q \otimes_A M \to Q \otimes_A M'$
is injective. Since localization is exact we see that
$(S')^{-1}(Q \otimes_A M) \to (S')^{-1}(Q \otimes_A M')$ is injective.
As $(S')^{-1}(Q \otimes_A M) = Q \otimes_A (S')^{-1}M$ and similarly for $M'$
we see that
$Q \otimes_A (S')^{-1}M \to Q \otimes_A (S')^{-1}M'$ is injective, hence
$(S')^{-1}M \to (S')^{-1}M'$ is universally injective as a map of
$A$-modules. This proves the first part of (1).
To see (2) we can use the following two facts: (a) if $Q$ is an
$S^{-1}A$-module, then $Q \otimes_A S^{-1}A = Q$, i.e., tensoring
with $Q$ over $A$ is the same thing as tensoring with $Q$ over $S^{-1}A$,
(b) if $M$ is any $A$-module on which the elements of $S$ are invertible,
then $M \otimes_A Q = M \otimes_{S^{-1}A} S^{-1}Q$.
Part (2) follows from this immediately.
\end{proof}

\begin{lemma}
\label{lemma-check-universally-injective-into-flat}
Let $R$ be a ring and let $M \to M'$ be a map of $R$-modules.
If $M'$ is flat, then $M \to M'$ is universally injective if
and only if $M/IM \to M'/IM'$ is injective for every finitely
generated ideal $I$ of $R$.
\end{lemma}

\begin{proof}
It suffices to show that $M \otimes_R Q \to M' \otimes_R Q$ is
injective for every finite $R$-module $Q$, see
Theorem \ref{theorem-universally-exact-criteria}.
Then $Q$ has a finite filtration
$0 = Q_0 \subset Q_1 \subset \ldots \subset Q_n = Q$
by submodules whose subquotients
are isomorphic to cyclic modules $R/I_i$, see
Lemma \ref{lemma-trivial-filter-finite-module}.
Since $M'$ is flat, we obtain a filtration
$$
\xymatrix{
M \otimes Q_1 \ar[r] \ar[d] &
M \otimes Q_2 \ar[r] \ar[d] &
\ldots \ar[r] &
M \otimes Q \ar[d] \\
M' \otimes Q_1 \ar@{^{(}->}[r] &
M' \otimes Q_2 \ar@{^{(}->}[r] &
\ldots \ar@{^{(}->}[r] &
M' \otimes Q
}
$$
of $M' \otimes_R Q$ by submodules $M' \otimes_R Q_i$ whose successive
quotients are $M' \otimes_R R/I_i = M'/I_iM'$. A simple induction argument
shows that it suffices to check $M/I_i M \to M'/I_i M'$ is injective.
Note that the collection of finitely generated ideals $I'_i \subset I_i$
is a directed set. Thus $M/I_iM = \colim M/I'_iM$ is a filtered
colimit, similarly for $M'$, the maps $M/I'_iM \to M'/I'_i M'$ are
injective by assumption, and since filtered colimits are exact
(Lemma \ref{lemma-directed-colimit-exact}) we conclude.
\end{proof}

\begin{lemma}
\label{lemma-base-change-along-universally-injective-ring-map}
Let $R \to S$ be a ring map which is universally injective
as a map of $R$-modules. Then the functor $M \mapsto M \otimes_R S$
on $R$-modules reflects injections, surjections, and isomorphisms.
\end{lemma}

\begin{proof}
Let $M \to N$ be a map of $R$-modules with kernel $K$ and cokernel $Q$.
If $M \otimes_R S \to N \otimes_R S$ is injective, then
the image of $K \otimes_R S \to M \otimes_R S$ is zero.
Since $K \subset K \otimes_R S$ and $M \subset M \otimes_R S$
we conclude that $K = 0$.
On the other hand, if $M \otimes_R S \to N \otimes_R S$ is surjective,
then $Q \otimes_R S$ is zero (by right exactness of tensor product).
Hence $Q \subset Q \otimes_R S$ is zero too.
\end{proof}

\begin{lemma}
\label{lemma-faithfully-flat-universally-injective}
Let $R \to S$ be a faithfully flat ring map.
Then $R \to S$ is universally injective as a map of $R$-modules.
In particular $R \cap IS = I$ for any ideal $I \subset R$.
\end{lemma}

\begin{proof}
Let $N$ be an $R$-module. We have to show that $N \to N \otimes_R S$ is
injective. As $S$ is faithfully flat as an $R$-module, it suffices to prove
this after tensoring with $S$. Hence it suffices to show that
$N \otimes_R S \to N \otimes_R S \otimes_R S$,
$n \otimes s \mapsto n \otimes 1 \otimes s$ is injective. This is true
because there is a retraction, namely,
$n \otimes s \otimes s' \mapsto n \otimes ss'$.
\end{proof}




